\documentclass[11pt]{article}
\usepackage{icmlko}
\begin{document}

\title{고르개는 없다 --- 세 시계가 저마다 다른 $1$ 등을 가리키고 맞은 것은 하나였다}
\author{월드모델 실험실}
\date{노트 501 · 2026년 8월 3일}
\maketitle

\begin{abstract}
노트 500 이 남긴 마지막 가설 --- \textbf{안쪽 검증은 유보에 붙어 있어야
한다.} 안쪽 검증을 $2024\!\sim\!2025$ 한 해로 좁혀 같은 $21$ 짝을 다시
쟀다.

\textbf{절반은 맞았다.} 인접할수록 순위가 선다 --- $2023\!\sim\!2025$
기준으로 $2022\!\sim\!2024$ 는 스피어만 $-0.320$ 인데
$2024\!\sim\!2025$ 는 $\mathbf{+0.728}$ 이다.

\textbf{그런데 $1$ 등의 정체는 안 선다.}

\begin{center}
\begin{tabular}{llrr}
\hline
안쪽 검증 & 그 시기의 $1$ 등 & 안쪽 판 차 & 유보 판 차 \\
\hline
$2022\!\sim\!2024$ & 모바일 media\_push & $+0.0035$ & (안 잼) \\
$2023\!\sim\!2025$ & 애니 트렌드 & $+0.0042$ & $\mathbf{+0.0046}$ \\
$2024\!\sim\!2025$ & 모바일 태그 & $+0.0097$ & $\mathbf{-0.0027}$ \\
\hline
\end{tabular}
\end{center}

\textbf{세 시기의 $1$ 등이 셋 다 다르고, 제일 방어 가능한 고르개(제일
인접 · 안쪽 판 · $1$ 등)가 뽑은 짝은 유보에서 진다}($-0.0027$ · $0/12$).

\textbf{$1/3$ 이다. 우연과 구별이 안 된다. 갈래를 닫는다.}
\end{abstract}

\begin{figure}[t]
\centering
\includegraphics[width=\textwidth]{figs/f1.pdf}
\caption{\textbf{가} --- 인접성과 순위 안정성. \textbf{나} --- 세 시기의
$1$ 등. \textbf{다} --- 유보 관문.}
\caption*{}
\end{figure}

\section{노트 498 채택의 지위}

\begin{center}
\begin{tabular}{lrr}
\hline
팔 & 판 차 & 부호 \\
\hline
모바일 태그만 & $-0.0027$ & $0/12$ \\
\textbf{애니 트렌드(노트 498)} & $\mathbf{+0.0046}$ & $\mathbf{12/12}$ \\
애니 $+$ 모바일 태그 & $+0.0016$ & $12/12$ \\
\hline
\end{tabular}
\end{center}

더하기는 성립한다($+0.0046 - 0.0027 = +0.0019$ 대 실제 $+0.0016$).

\textbf{채택은 그대로 둔다} --- 유보 관문에서 자 다섯이 전부 양수였고
($10/10$), 그 값은 어떻게 고랐는지와 무관하다. \textbf{그러나 이제
분명하다 --- 그 짝을 짚은 것은 고르개가 아니라 운이었다.} 두 해 동안의
다른 모든 채택과 같은 지위다: \emph{후보를 만들고 관문에 태워 통과한 것.}

\begin{quote}
\textbf{조항} --- \textbf{고르개가 있다고 말하려면 서로 다른 안쪽 시기
셋에서 같은 $1$ 등이 나와야 한다.} 이 자료에서는 안 나온다. 한 시기의
$1$ 등이 유보에서 이겼다는 사실은 \textbf{고르개의 증거가 아니라 그
후보의 증거}다.
\end{quote}

\section{그래도 남는 것 --- 인접성은 진짜다}

$-0.320 \to +0.728$ 은 큰 차이다. \textbf{안쪽 검증을 유보에 붙이면
순위가 선다.} 다만 그 순위는 \emph{안쪽 시기끼리} 서는 것이지 유보와
서는 것이 아니다.

읽으면 --- 안쪽 판 차는 \textbf{그 시기의 진짜 성질}을 재고 있다(그래서
가까운 시기끼리 닮는다). 다만 \textbf{$2025$ 년의 성질은 $2024$ 년에도
아직 없다.} 노트 493 이 애니 트렌드 덮음이 $2025$ 에 $55.6\% \to 96.5\%$
로 뛰는 것을 셌는데, 그런 일이 안쪽에는 안 들어 있다.

\section{다음 길}

안쪽 검증으로 고르는 길은 닫혔다. 남은 길은 둘이다.

\begin{enumerate}
\item \textbf{유보를 나눠 쓴다} --- $2025$ 상반기로 고르고 하반기로
채점. 채점 집합이 바뀌므로 두 해치 기록과 못 견준다(노트 $82 \cdot 90$
분모 규칙). \textbf{값이 큰 갈래에만 쓴다.}
\item \textbf{후보를 하나씩 관문에 태운다} --- 이 실험실이 두 해 동안 한
바로 그것. 노트 497 이 이미 도메인별 값을 다 재 놨다(웹툰
\texttt{target\_breadth} $+0.0130$ · 모바일 $+0.0157$ · 시장팝업
$+0.0553$ · 게임 $+0.0208$). \textbf{판 무게를 곱하면 웹툰이 제일 크다.}
\end{enumerate}

\textbf{②로 간다.} 고르개를 찾는 대신 관문을 여러 번 쓴다 --- 다만
자 다섯을 전부 걸어서, 집 밖 두 짝이 독립 검증이 되게 한다.

\section*{남은 것}

\begin{itemize}
\item \textbf{웹툰 \texttt{target\_breadth}} --- 기대 이득 $+0.0027$ 로
제일 크다. 다음 관문.
\item \textbf{남은 여덟 축의 도메인별 값} --- 아직 안 쟀다. 노트 497 은
셋만 쟀다.
\item \textbf{인접성 발견을 어디에 쓰나} --- 순위가 시기끼리 선다면,
\emph{시기별로 축을 다르게 쓰는} 모형이 가능한가. 지금 구조에는 자리가
없다.
\end{itemize}

\end{document}
